\section{Conclusion}

This paper studies tool-using LLM agents from the perspective of explicit orchestration. Rather than relying entirely on fixed workflows or unconstrained prompt-driven reasoning, we formulate agent control as a utility-guided decision process over actions such as responding, retrieving, calling tools, verifying, and stopping. The resulting framework makes it possible to analyze how estimated gain, cost proxies, uncertainty, and redundancy shape the quality--cost trade-off of agent execution.

Our experiments show that explicit utility design provides a practical and defensible way to study agent behavior, even when the underlying signals are heuristic. Additional analyses of cost definitions, workflow fairness, and redundancy control further clarify what the policy is doing and what it is not doing. Overall, the main value of this work is not a claim of universal dominance, but a structured step toward more controllable, analyzable, and cost-aware LLM agents.
We believe this perspective is particularly relevant for budget-sensitive and context-limited deployments, where multi-step agent behavior must be effective not only in quality, but also in cost and control.